\section{Redes Definidas por Software (SDN)}

As Redes Definidas por Software (SDN, Software-Defined Networking) apresentam uma abordagem na qual o plano de controle é separado do plano de dados da rede. Dessa forma, os dispositivos responsáveis pelo encaminhamento dos dados, como switches e roteadores, podem receber regras definidas por uma camada de controle, permitindo que o comportamento da rede seja gerenciado de maneira programável \cite{bera2017software}.

A arquitetura SDN pode ser organizada em três camadas principais: infraestrutura, controle e aplicação. A camada de infraestrutura é composta pelos dispositivos responsáveis pelo encaminhamento dos dados, enquanto a camada de controle concentra a lógica responsável pelo gerenciamento da rede. Já a camada de aplicação permite a implementação de aplicações e serviços que utilizam as informações fornecidas pelo controlador. A comunicação entre essas camadas ocorre por meio de interfaces que possibilitam a troca de informações e a definição das regras aplicadas à infraestrutura \cite{bera2017software}.

A possibilidade de programar o comportamento da rede proporciona maior flexibilidade em comparação às arquiteturas tradicionais, permitindo que suas configurações sejam modificadas de acordo com diferentes requisitos \cite{bera2017software}. Essa característica contribui para o fatiamento de redes, no qual diferentes fatias podem exigir configurações específicas e adaptações na utilização dos recursos disponíveis \cite{gomes2020strategies}.